\documentclass[12pt]{article}

% =========================
% 0) Metadata (EDIT ME)
% =========================
\newcommand{\CourseCode}{MATH 431}
\newcommand{\CourseTitle}{Introduction to Probability}
\newcommand{\CourseSection}{LEC 002} % change to 001 or 003 if needed
\newcommand{\Semester}{Spring 2026}
\newcommand{\Instructor}{Sarah Tammen} % LEC 001: Ander Aguirre | LEC 002: Sarah Tammen | LEC 003: Nicholas Derr
\newcommand{\MeetingInfo}{MWF 9:55--10:45, 6102 Sewell Social Sciences} % LEC 001 default; update if in 002/003
\newcommand{\HomeworkNumber}{10}
\newcommand{\YourName}{Alejandro De La Torre}
\newcommand{\DueDate}{April 17, 2026} % or e.g. February 2, 2026

% =========================
% 1) Core math tools
% =========================
\usepackage{amsmath, amssymb, amsthm}

% =========================
% 2) Lists and spacing
% =========================
\usepackage{enumitem}
\setlist[enumerate,1]{label=\textbf{(\alph*)}, leftmargin=2em, itemsep=0.25em}
\setlist[itemize]{leftmargin=1.5em, itemsep=0.25em}
\setlength{\parskip}{0.35em}
\setlength{\parindent}{0pt}

% =========================
% 3) Page geometry + header/footer
% =========================
\usepackage{geometry}
\geometry{margin=1in}

\usepackage{fancyhdr}
\pagestyle{fancy}
\fancyhf{}
\lhead{\CourseCode\ --- \CourseSection, \Semester}
\rhead{Homework \HomeworkNumber}
\cfoot{\thepage}
\renewcommand{\headrulewidth}{0.4pt}
\setlength{\headheight}{15pt}
\addtolength{\topmargin}{-2pt}

% =========================
% 4) Links (subtle)
% =========================
\usepackage[hidelinks]{hyperref}
\setcounter{tocdepth}{1}

% =========================
% 5) Figures (optional)
% =========================
\usepackage{graphicx}
\graphicspath{{images/}}

% =========================
% 6) Lightweight boxes (optional)
% =========================
\usepackage{mdframed}
\newmdenv[
  skipabove=0.6\baselineskip,
  skipbelow=0.6\baselineskip,
  linewidth=0.6pt,
  roundcorner=4pt,
  innerleftmargin=10pt,
  innerrightmargin=10pt,
  innertopmargin=8pt,
  innerbottommargin=8pt
]{boxnote}

% =========================
% 7) Probability / notation macros
% =========================
\newcommand{\R}{\mathbb{R}}
\newcommand{\N}{\mathbb{N}}
\newcommand{\Z}{\mathbb{Z}}

\newcommand{\OmegaSpace}{\Omega}
\newcommand{\FField}{\mathcal{F}}
% Probability measure in case it is relevant or want to distinguish it from P(A)
%\newcommand{\PMeasure}{\mathbb{P}} 
\newcommand{\E}{\mathbb{E}}

\newcommand{\given}{\,\middle|\,}
\newcommand{\ind}{\mathbf{1}}

% Common “defined as” symbol
\newcommand{\defeq}{\vcentcolon=}

% =========================
% 8) Problem command (TOC + PDF bookmarks)
% Usage:
%   \problem{<label>}
%   \problem[Short title]{<label>}
% Example:
%   \problem{1.4}
%   \problem[Sample space]{1.4}
% =========================
\makeatletter
\newcommand{\problem}[2][]{%
  \def\prob@title{Problem #2\if\relax\detokenize{#1}\relax\else\ (\textit{#1})\fi}%
  \phantomsection%
  \pdfbookmark[1]{\prob@title}{prob#2}%
  \addcontentsline{toc}{section}{\prob@title}%
  \section*{\prob@title}%
  \vspace{-0.25em}\hrule\vspace{0.8em}%
}
\makeatother

% =========================
% 9) Solution environment
% =========================
\newenvironment{solution}{%
  \noindent\textbf{Solution.}\ \ignorespaces
}{\par}

% Optional scratch / planning block
\newenvironment{scratch}{%
  \begin{boxnote}\textbf{Scratch / Setup.}\ \small
}{\end{boxnote}}

% Final answer command: expects math only (no $$ or \[ \])
\newcommand{\finalanswer}[1]{%
  \par\medskip
  \noindent\textbf{Final:}\ \fbox{$#1$}\par
} % AGAIN: MATH ONLY INSIDE THE BOX; if word answer, use \text{} inside or use \fbox{TEXT} or \framebox

% =========================
% 10) Title block
% =========================
\title{Homework \HomeworkNumber}
\author{\YourName \\
\CourseCode\ --- \CourseTitle \\
\CourseSection \\
Instructor: \Instructor}
\date{Due: \DueDate}

\begin{document}
\maketitle

% \section*{Collaboration / Resources}
% (Optional) Write “I worked alone” or list collaborators/resources.

\tableofcontents
\bigskip
\hrule
\bigskip
\newpage

% ============================================================
% Problems (duplicate blocks as needed)
% ============================================================


\problem{7.2}
Let $X$ and $Y$ be independent Bernoulli random variables with parameters $p$ and $r$, respectively. Find the distribution of $X+Y$.

\begin{solution}

\begin{scratch}
Let
\[
S=X+Y.
\]
Since $X,Y\in\{0,1\}$, the only possible values of $S$ are $0,1,2$.

The relevant combinations are:
\[
\begin{aligned}
S=0 &\iff (X,Y)=(0,0),\\
S=1 &\iff (X,Y)\in\{(1,0),(0,1)\},\\
S=2 &\iff (X,Y)=(1,1).
\end{aligned}
\]
Because $X$ and $Y$ are independent Bernoulli random variables,
\[
P(X=1)=p,\quad P(X=0)=1-p,\quad P(Y=1)=r,\quad P(Y=0)=1-r.
\]
\end{scratch}

We are computing the probability mass function of $S=X+Y$.

Since $S$ can take only the values $0,1,2$, we compute these probabilities directly.
First,
\[
P(S=0)=P(X=0,Y=0)=P(X=0)P(Y=0)=(1-p)(1-r).
\]
Next,
\[
P(S=2)=P(X=1,Y=1)=P(X=1)P(Y=1)=pr.
\]
Finally,
\begin{align*}
P(S=1)
&=P(X=1,Y=0)+P(X=0,Y=1)\\
&=P(X=1)P(Y=0)+P(X=0)P(Y=1)\\
&=p(1-r)+(1-p)r\\
&=p+r-2pr.
\end{align*}

Therefore the pmf of $X+Y$ is
\[
p_S(s)=P(S=s)=
\begin{cases}
(1-p)(1-r), & s=0,\\[4pt]
p+r-2pr, & s=1,\\[4pt]
pr, & s=2,\\[4pt]
0, & \text{otherwise}.
\end{cases}
\]
This is a binomial distribution only in the special case $p=r$.

\finalanswer{P(X+Y=0)=(1-p)(1-r),\quad P(X+Y=1)=p+r-2pr}
\finalanswer{P(X+Y=2)=pr}

\end{solution}

\newpage

\problem{7.16}
Let $X$ be a Poisson random variable with parameter $\lambda$, and $Y$ an independent Bernoulli random variable with parameter $p$. Find the probability mass function of $X+Y$.

\begin{solution}

\begin{scratch}
Let
\[
S=X+Y.
\]
Since $X\in\{0,1,2,\dots\}$ and $Y\in\{0,1\}$, the support of $S$ is also
\[
\{0,1,2,\dots\}.
\]

The cleanest computation is to condition on $Y$:
\[
P(S=n)=P(Y=0)P(X=n)+P(Y=1)P(X=n-1),
\]
with the convention that $P(X=-1)=0$.

Also,
\[
P(X=k)=e^{-\lambda}\frac{\lambda^k}{k!},\qquad k=0,1,2,\dots
\]
and
\[
P(Y=0)=1-p,\qquad P(Y=1)=p.
\]
\end{scratch}

We want the probability mass function of $S=X+Y$.

For $n\ge 0$,
\[
P(S=n)=P(Y=0)P(X=n)+P(Y=1)P(X=n-1).
\]
If $n=0$, the second term is zero, so
\[
P(S=0)=(1-p)P(X=0)=(1-p)e^{-\lambda}.
\]

If $n\ge 1$, then
\begin{align*}
P(S=n)
&=(1-p)e^{-\lambda}\frac{\lambda^n}{n!}
+p\,e^{-\lambda}\frac{\lambda^{n-1}}{(n-1)!}\\
&=e^{-\lambda}\frac{\lambda^{n-1}}{n!}\bigl((1-p)\lambda+np\bigr).
\end{align*}

Hence the pmf of $X+Y$ is
\[
p_{X+Y}(n)=
\begin{cases}
(1-p)e^{-\lambda}, & n=0,\\[6pt]
e^{-\lambda}\dfrac{\lambda^{n-1}}{n!}\bigl((1-p)\lambda+np\bigr), & n=1,2,3,\dots
\end{cases}
\]

\finalanswer{p_{X+Y}(0)=(1-p)e^{-\lambda},\quad p_{X+Y}(n)=e^{-\lambda}\frac{\lambda^{n-1}}{n!}\bigl((1-p)\lambda+np\bigr)\ \text{for }n\ge 1}

\end{solution}

\newpage

\problem{7.20}
Let $X$ have density $f_X(x)=2x$ for $0<x<1$ and 0 outside of $(0,1)$, and let $Y$ be uniform on the interval $(1,2)$. Assume $X$ and $Y$ independent.
\begin{enumerate}[label=\textbf{(\alph*)}]
  \setcounter{enumi}{1}
  \item Find the density function of $X+Y$.
\end{enumerate}

\begin{solution}

\begin{scratch}
Let
\[
T=X+Y.
\]
Since $X\in(0,1)$ and $Y\in(1,2)$, we must have
\[
T\in(1,3).
\]

By convolution,
\[
f_T(t)=\int_{-\infty}^{\infty} f_X(x)f_Y(t-x)\,dx.
\]
Here
\[
f_X(x)=
\begin{cases}
2x, & 0<x<1,\\
0, & \text{otherwise},
\end{cases}
\qquad
f_Y(y)=
\begin{cases}
1, & 1<y<2,\\
0, & \text{otherwise}.
\end{cases}
\]

The integrand is nonzero exactly when
\[
0<x<1
\quad\text{and}\quad
1<t-x<2,
\]
that is,
\[
\max(0,t-2)<x<\min(1,t-1).
\]
This leads to two cases:
\[
1<t<2 \quad\text{and}\quad 2\le t<3.
\]
\end{scratch}

We are finding the density of $T=X+Y$.

For any real $t$,
\[
f_T(t)=\int_{-\infty}^{\infty} f_X(x)f_Y(t-x)\,dx.
\]
If $t\notin(1,3)$, there is no overlap of the supports, so $f_T(t)=0$.

If $1<t<2$, then $\max(0,t-2)=0$ and $\min(1,t-1)=t-1$, so
\[
f_T(t)=\int_0^{t-1}2x\,dx=(t-1)^2.
\]

If $2\le t<3$, then $\max(0,t-2)=t-2$ and $\min(1,t-1)=1$, so
\[
f_T(t)=\int_{t-2}^{1}2x\,dx=1-(t-2)^2.
\]

Therefore
\[
f_{X+Y}(t)=
\begin{cases}
(t-1)^2, & 1<t<2,\\[4pt]
1-(t-2)^2, & 2\le t<3,\\[4pt]
0, & \text{otherwise}.
\end{cases}
\]
As a quick check,
\[
\int_1^2 (t-1)^2\,dt+\int_2^3 \bigl(1-(t-2)^2\bigr)\,dt
=\frac13+\frac23=1,
\]
so this is a valid density.

\finalanswer{f_{X+Y}(t)=\begin{cases}(t-1)^2,&1<t<2\\ 1-(t-2)^2,&2\le t<3\\ 0,&\text{otherwise}\end{cases}}

\end{solution}

\newpage

\problem{7.24}
Suppose that $X$, $Y$, and $Z$ are independent mean zero normal random variables. Find $P(X+2Y-3Z>0)$.

\begin{solution}

\begin{scratch}
Let
\[
W=X+2Y-3Z.
\]
Because $X,Y,Z$ are independent normal random variables, any linear combination
of them is also normal. So $W$ is normal.

Its mean is
\[
E[W]=E[X]+2E[Y]-3E[Z]=0.
\]
If we write
\[
\operatorname{Var}(X)=\sigma_X^2,\quad \operatorname{Var}(Y)=\sigma_Y^2,\quad \operatorname{Var}(Z)=\sigma_Z^2,
\]
then independence gives
\[
\operatorname{Var}(W)=\sigma_X^2+4\sigma_Y^2+9\sigma_Z^2.
\]
So $W$ is a mean-zero normal random variable, hence its distribution is symmetric
about $0$.
\end{scratch}

We are asked to compute the probability
\[
P(X+2Y-3Z>0)=P(W>0),
\]
where
\[
W=X+2Y-3Z.
\]

Since $W$ is a linear combination of independent normal random variables, it is
itself normal. Its mean is
\[
E[W]=0,
\]
and its variance is
\[
\operatorname{Var}(W)=\operatorname{Var}(X)+4\operatorname{Var}(Y)+9\operatorname{Var}(Z).
\]
Therefore $W$ is a centered normal random variable. A mean-zero normal
distribution is symmetric about $0$, so
\[
P(W>0)=\frac12.
\]

\finalanswer{P(X+2Y-3Z>0)=\frac12}

\end{solution}

\newpage

\problem{7.28}
Let $X_1,X_2,X_3$ be independent $\operatorname{Exp}(\lambda)$ distributed random variables. Find the probability that $X_1<X_2<X_3$.

\begin{solution}

\begin{scratch}
The variables $X_1,X_2,X_3$ are i.i.d. and continuous.

Because they are continuous,
\[
P(X_i=X_j)=0\qquad \text{for } i\ne j,
\]
so ties occur with probability $0$.

There are $3!=6$ strict orderings:
\[
\begin{aligned}
&X_1<X_2<X_3,\qquad X_1<X_3<X_2,\qquad X_2<X_1<X_3,\\
&X_2<X_3<X_1,\qquad X_3<X_1<X_2,\qquad X_3<X_2<X_1.
\end{aligned}
\]
By exchangeability of i.i.d. random variables, these six orderings are equally likely.
\end{scratch}

Since $X_1,X_2,X_3$ are independent and identically distributed continuous random
variables, each strict ordering of the three variables has the same probability.
Also, because the distribution is continuous, ties have probability zero. Hence the
six events
\[
\begin{aligned}
&X_1<X_2<X_3,\qquad X_1<X_3<X_2,\qquad X_2<X_1<X_3,\\
&X_2<X_3<X_1,\qquad X_3<X_1<X_2,\qquad X_3<X_2<X_1
\end{aligned}
\]
are disjoint, equally likely, and their union has probability $1$.

Therefore each one has probability
\[
\frac{1}{6}.
\]

\finalanswer{P(X_1<X_2<X_3)=\frac16}

\end{solution}

\newpage

\problem{7.30}
We deal five cards, one by one, from a standard deck of 52.
(Dealing cards from a deck means sampling without replacement.)
\begin{enumerate}[label=\textbf{(\alph*)}]
  \item Find the probability that the second card is an ace and the fourth card is a king.
  \item Find the probability that the first and the fifth cards are both spades.
  \setcounter{enumi}{2}
  \item Find the conditional probability that the second card is a king given that the last two cards are both aces.
\end{enumerate}

\begin{solution}

\begin{scratch}
We deal five cards from a uniformly shuffled deck, so all ordered 5-card sequences
are equally likely.

Useful facts:
\begin{itemize}
  \item There are $4$ aces, $4$ kings, and $13$ spades.
  \item By symmetry/exchangeability, specified positions behave like any other
  specified positions.
  \item Conditioning on the last two cards being aces leaves $50$ cards for the
  first three positions, still including all $4$ kings.
\end{itemize}
\end{scratch}

We compute each part separately.

\begin{enumerate}[label=\textbf{(\alph*)}]
  \item We need the probability that the second card is an ace and the fourth card
  is a king.

  By exchangeability of the positions in a random ordering of the deck, this is the
  same as the probability that the first card is an ace and the second card is a king.
  Therefore
  \[
  P(\text{2nd ace and 4th king})
  =\frac{4}{52}\cdot\frac{4}{51}
  =\frac{16}{2652}
  =\frac{4}{663}.
  \]

  \item We need the probability that the first and fifth cards are both spades.

  Again by exchangeability, this is the same as the probability that the first two
  cards are both spades. Thus
  \[
  P(\text{1st and 5th spades})
  =\frac{13}{52}\cdot\frac{12}{51}
  =\frac{1}{17}.
  \]

  \item We need
  \[
  P(\text{2nd card is a king}\mid \text{last two cards are both aces}).
  \]

  After conditioning on the last two cards being aces, the first three positions are
  filled by an ordered sample from the remaining $50$ cards. None of the aces is a
  king, so all $4$ kings are still among those $50$ cards. By symmetry, each of the
  first three positions is equally likely to contain any one of the remaining $50$
  cards. Hence
  \[
  P(\text{2nd card is a king}\mid \text{last two cards are both aces})
  =\frac{4}{50}
  =\frac{2}{25}.
  \]
\end{enumerate}

\finalanswer{\text{(a)}\ \frac{4}{663}}
\finalanswer{\text{(b)}\ \frac{1}{17}}
\finalanswer{\text{(c)}\ \frac{2}{25}}

\end{solution}


\end{document}
